\section{Conclusion and Future Directions}
\label{sec:conclusion}

This paper presented \name, an open AI-powered command-line agent for software engineering, and documented the architectural decisions, design trade-offs, and lessons learned from building a production-ready system. Key contributions include a compound AI system architecture~\cite{zaharia2024compound} with multi-model routing (\Cref{sec:multi_model}), an extended ReAct pipeline with thinking and critique phases (\Cref{sec:react_executor}), Adaptive Context Compaction (\Cref{sec:adaptive_compaction}), event-driven system reminders (\Cref{sec:system_reminders}), an experience-driven memory pipeline, lazy tool discovery via MCP (\Cref{sec:mcp}), and a dual-interface abstraction.

The central architectural insight is that per-workflow LLM binding (\Cref{sec:multi_model}) yields model-agnosticity: adapting to new models requires only configuration changes, not code changes. Schema-level safety enforcement (\Cref{sec:agent_core}) proved more robust than runtime permission checks, as removing write tools from the planning agent's schema eliminates an entire class of bypass attempts. Conditional prompt composition (\Cref{sec:prompt_composition}) reduced overhead by excluding irrelevant instructions while preserving comprehensive guidance when needed. On the context engineering side, managing the finite context window emerged as a first-class engineering concern rather than a secondary optimization. Adaptive Context Compaction (\Cref{sec:adaptive_compaction}), which transitions observations through active, faded, and archived states, reduced peak context consumption by approximately 54\% and often eliminated the need for emergency summarization. The three-tier context architecture (static system prompt, dynamic just-in-time reminders (\Cref{sec:system_reminders}), and long-horizon persistence) addressed the attention-decay problem where agents reliably violated instructions after 30+ tool calls. The Agentic Context Engineering memory pipeline enables the agent to learn from tool outcomes within and across sessions without hardcoding strategies into prompts.

The cross-cutting design tensions and transferable lessons that emerged from this work are synthesized in \Cref{sec:discussion}, which examines context pressure as a first-class constraint, behavioral steering over long horizons, safety through architectural enforcement rather than runtime checks, tool design that absorbs LLM imprecision, and resource bounding strategies for long-running sessions.

Several capabilities previously identified as future work, such as per-step undo via shadow git snapshots (\Cref{sec:undo}), has been implemented during continued development, demonstrating the extensibility of the layered architecture.

Several promising research directions emerge from the challenges identified in this work:

\begin{packeditemize}
\item \textbf{Quantitative evaluation on established benchmarks.} This paper documents architectural decisions and design rationale but lacks systematic quantitative evaluation. Benchmarking on SWE-bench~\cite{yang2024swe}, Terminal-Bench~\cite{terminal_bench}, and LongCLI-Bench~\cite{longcli_bench}, part of the broader evaluation ecosystem surveyed in \Cref{sec:related_work}, would validate the architecture against established baselines and identify specific improvement areas. Terminal-Bench's finding that frontier agents resolve fewer than 65\% of tasks, and LongCLI-Bench's observation that pass rates remain below 20\% for long-horizon tasks, suggest substantial room for improvement in context management and multi-step reasoning.

\item \textbf{Adaptive resource allocation.} Current parameters - the 70\% compaction threshold, 3 nudge attempts, 6 thinking depth levels - are globally fixed. Adaptive approaches that dynamically adjust based on task complexity, current context pressure, and error history could optimize the cost--quality--latency triangle per task rather than relying on one-size-fits-all constants. For instance, a simple debugging task should skip the thinking phase entirely, while a complex architectural refactoring may benefit from deeper deliberation and a more conservative compaction strategy.

\item \textbf{Scaling the memory pipeline.} The ACE playbook currently operates per-project with effectiveness scoring and semantic retrieval. Cross-project knowledge transfer - where lessons learned in one repository inform behavior in similar projects - hierarchical bullet organization that separates general programming heuristics from project-specific conventions, and active learning that selectively requests user feedback on uncertain bullets could substantially improve the agent's ability to accumulate and apply experience over time.

\item \textbf{Structured code representations for memory.} The current memory pipeline stores experience as flat natural-language bullets. Richer representations, such as code dependency graphs that capture relationships between modules, call graphs that track function interactions, and project-level ontologies that encode domain concepts, could enable more precise retrieval and reasoning. Combining graph-structured code understanding with long-term persistent memory that spans not just sessions but entire project lifecycles would allow the agent to build deep, evolving models of codebases rather than accumulating isolated observations.

\item \textbf{Multi-agent coordination beyond hierarchical delegation.} Current subagents execute independently under main-agent coordination, communicating only through completion markers. Richer coordination patterns - peer-to-peer communication between subagents, shared blackboard architectures for collaborative problem-solving, and negotiation protocols for resolving conflicting tool outcomes - could enable more sophisticated workflows such as concurrent code review and implementation, or parallel exploration with result synthesis.

\item \textbf{Learned system reminder optimization.} The 24-template reminder catalog and its injection timing were manually engineered based on observed failure modes. Automated discovery of effective reminder patterns - through reinforcement learning on attention-decay metrics, learned injection timing based on conversation dynamics, or adaptive template selection conditioned on agent state - could improve nudge effectiveness while reducing the manual engineering burden of designing and maintaining reminder templates.

\item \textbf{Hybrid CLI--IDE integration.} The dual-interface architecture demonstrates that a shared callback protocol can serve fundamentally different frontends. Extending this to IDE plugins - where the same agent logic powers both terminal workflows and rich editor integration - would serve users who want visual affordances (inline diffs, symbol navigation, test result overlays) alongside terminal autonomy, without duplicating agent logic across environments.
\end{packeditemize}

Building effective agentic coding systems requires navigating competing concerns: capability versus complexity, autonomy versus safety, generality versus token efficiency. The design space is rich with trade-offs where no single choice dominates. We hope that the architectural patterns, engineering lessons, and candid reflections on what worked and what did not, documented in this paper, help future builders of agentic systems make more informed decisions as they navigate these choices.
